\documentclass[../main.tex]{subfiles}

\begin{document}

\chapter{Metode Pembuktian Matematis (Langsung, Kontradiksi, Kontraposisi)}

\begin{subcpmk}
  \item Mahasiswa mampu menerapkan beragam kaidah dan spesifikasi dari metode pembuktian matematis terhadap analisis fungsi kalkulasi logis dan pembuktian algoritma komputasi awal.
\end{subcpmk}

\section{Urgensi Pembuktian (Proofs) dalam Pengembangan Perangkat Lunak}

Mengapa seorang perancang algoritma Teknik Informatika wajib terbiasa dengan metode pembuktian matematis? Sebab validasi empiris komputasi kompilator / \text{Testing Kode Prosedural} belumlah cukup menjamin kekonsistenan perangkat lunak di masa depan. *Unit Test* dalam *software deployment* hanya menguji sejauh mana tebakan kondisi input yang programmer buat divaluasi *output* kebenarannya. Saat ruang input fungsi berjumlah milyaran kemungkinan iteratif, kita butuh "Bukti Formal (*Formal Verification*)" secara rasional dan aljabar logis bahwa sistem selamanya merespons iterasi parameter komputasional berdaya absolut sah (Selamanya tak bakal \textit{Crash} atau membiarkan malfungsi data komputasi bocor). Argumentasi \textit{mathematical proof}-lah piranti mumpuni jaminan matematis arsitektur ini.

Secara teknis, \textbf{Bukti (Proof)} dideskripsikan berupa urutan barisan klaim kebenaran (proposisi logikal absolut takterbantahkan $\to$ ke sebuah kebenaran baru $\to$ ditarik ke \textit{simpul kesimpulan fungsional}) yang secara solid merajarkan argumen algoritmanya sahih (valid). Komponen penyangga kebenaran berantai ini bersumber atas deduksi, dalil Tautologi ekivalensi, atau dari Aksioma yang sudah universal faktanya (Misal: "Sistem paritas aritmatika Genap bisa diejawantahkan ke perumusan $2k$ bil bulat").

\section{Tipologi Teknik Pembuktian Dasar Formal}
Buku ini mem-bina *student* menapaki alur utama 3 tipe teknik bukti: (Pembuktian Langsung), dan cara bukti tak-langsung (Via Kontra-posisi dan deduksi iterasi Kontradiksi logikal).

\subsection{1. Bukti Langsung (Direct Proof)}
Pada formasi fungsi hipotesis kausalitas logikal $(A \rightarrow B)$. Strategi pembuktian analitikal *Direct* mendikte ke kita: Terimalah dan asumsikan \textit{input} atau anteseden subjeknya (Klausa komputasional $A$) adalah BENAR secara murni. Lalu modifikasi spesifikasi rumus, faktorisasi persamaan / \textit{logical string} secara berurut valid hingga rumus itu menyala-nyala berkesudahan berbukti di posisi terminal absolut konklusi output *B* juga TERBUKTI murni (True T).

\textbf{Contoh (Klaim: Jika \textit{n} adalah parameter nilai algoritma tipe \textit{integer} ganjil, maka output \textit{kuadrat fungsi numerik parameternya} selalu berspesifikasi bilangan komputatif berpredikat ganjil) :}
- Asumsikan input klausa A adalah Mutlak: "$n$ wajib bilangan Ganjil algoritmanya".
- Ciri variabel numerik iterasi integer ganjil rasional dirumuskan definisi universal menjadi formula numerik dasar ($n = 2k + 1$).
- Substitusikan ke bentuk klausa terminal fungsi komputasi algoritmanya \textit{(Output evaluatif $B$ yakni $n^2$)}.
- Formula fungsi algoritmis akhir outputnya = $(2k+1)^2 = 4k^2 + 4k + 1$.
- Sederhanakan faktorisasi parameter iteratif terminal fungsinya: tarik faktor persilangan komputatif menjadi bernilai $2(2k^2 + 2k) + 1$.
- Karena $(2k^2 + 2k)$ secara absolut dinilai juga pasti bernilai \textit{Integer (variabel bulat)} kita misal rumuskan sebagai variabel konstanta $m$, kesimpulan rumusan komputasinya berekonstruksi kembali di pakem universal fungsi = $2(m) + 1$.
- **Kesimpulan Valid Absolut**: Terbukti mutlak format formula $(2m + 1)$ selalu mendeskripsikan secara sah tipe sirkulasi var \textit{Ganjil}. Terbukti \textit{output $B$} benar.

\subsection{2. Bukti Kontraposisi (Proof by Contraposition/Contrapositive)}
Teknik validasi ini adalah "jalan belakang" sistematis. Kita menggunakan Hukum Ekivalensi Logikal bahwasanya relasi fungsional parameter komputasi $p \rightarrow q$ sesungguhnya bernilai logik kembar mutlak ekuivalensinya dengan membalik posisi serta menyangkali klausa term (\textit{Tollens mode manipulation}): $\sim q \rightarrow \sim p$.
Metode: Mulai pengujian fungsional algoritmanya dengan mengandaikan Output Sistem komputasi Batal (B \textit{merupakan hal Salah / False}). Gunakan penjabaran deduktif sistematisnya menuju kesimpulan final bahwa Prasyarat var Input $A$ mestilah logikalnya ikut Batal (Salah / Negasi / \textit{False}). 

\textbf{Contoh (Klaim: Pada komputasi variabel numerik integer \textit{n}, jika resultan komputasi fungsi $3n + 2$ merender variabel balikan/output sirkulasi ganjil, maka logikal kepastian var \textit{n}-nya di pakem Input wajib ganjil mutlak).}
- Ini format kausal parameter komputasional: ($3n+2$ GANJIL $\rightarrow$ \textit{n} GANJIL).
- Buktikan relasi tak logis dengan Kontraposisi: (Asumsi kebalikan \textit{n} TIDAK GANJIL [Artinya komputasi bernilai *n Genap*], dan buktikan \textit{output} \textit{Resultan $3n + 2$ TIDAK GANJIL}).
- Mulai dari negasi *output* hipotesis fungsi \textit{(yaitu n adalah berformat num komputatif Genap)} $\to n = 2k$.
- Hitung evaluasi terminal parameter rumusnya: $3(2k) + 2$ algoritmanya dikonversi logikal menjadi formula algebraik konstan $= 6k + 2$.
- Faktorkan parameternya rasional: format persamaan berwujud formula $2(3k + 1)$.
- Konsekuensi akhir parameter: karena var konstan berbentuk \textit{kali 2 alias [2m]}, pasti tipe struktur formula memvonis bilangannya bernilai fungsional iteratif *GENAP*. 
- Kesimpulan formal absolut \textit{software}: $\to$ Benar Negasi Input. (Jika \textit{n} bernilai Genap inputnya, logikal terbukti output bernilai predikatif Genap/Bukan Ganjil). Klaim original sistem algoritmiknya: VALiD diimplementasi nyata!

\subsection{3. Bukti Kontradiksi (Proof by Contradiction / \textit{Reductio ad Absurdum})}
Teknik validasi *Testing* \textit{Crash} ekstrem. Buktikan bahwa hal berkebalikan malah melahirkan Paradoks yang melanggar Aksioma universal logikal!
Kita menyanggah total rasional kausal dari Klaim komputatif $(A \rightarrow B)$ secara utuh dengan cara memaksa sistem mengasumsikan input syarat A diset bernilai benar namun logikal output variabel B dikondisikan salah \textit{False/Negasi} mutlak $(A \land \sim B)$.
Jika dari iterasi jalan perhitungan ini lantas secara formal tersandung konklusi \textit{error komputatif irrasional (contoh: kompilator mendapat nilai komputasional hitung bahwa $1 = 0$, atau hasil komputasi logis var $x$ bernilai paritas ganjil dan paritas angka genap SEKETIKA BERSAMAAN secara absolut [Sebuah kemustahilan sistemik kontradiksi paradox!!!])}, maka asumsi paradoks ini kita Coret/Tolak mentah-mentah!. Dan konklusi mutlak berbalik menyatakan secara tunggal hanya Klaim original Algoritma sistemlah yang VALID TERBUKTI dan SAH! \textit{Q.E.D (Quod Erat Demonstrandum - Which was to be demonstrated)}.

\begin{aktivitas}
  \item Diskusikan pada struktur blok tim komputasi logika: Kapan menurut tim logikal rasional bahwa metode pengandaian Contrapositive (\textit{Contraposition}) lebih gampang digunakan di fungsi algoritmik sistem numerik iteratif dibandingkan jalan membuktikan rumus fungsi var \textit{Direct Proof}?
\end{aktivitas}


\begin{latihan}
  \item Validasilah *code-rule komputasi spesifikasi* logika \textit{array loop} ini menggunakan \textit{Metode Bukti Contraposisi (Tak-langsung)} mutlak: "Jika logikal komputasional bilangan bulat \textit{var integer} $x^2$ diinspeksi bertipe kalkulasi genap, maka sudah \textit{pasti dan absolut mutlak} parameter fungsi var mandirinya \textit{x} dispesifikasi sebagai bilangan bulat bertipe genap rasional!".
\end{latihan}

\begin{asesmen}
Soal penalaran logika komputasi fungsional di portofolio kuis. Peniliaian ditelaah dari bagaimana tahapan algoritma penyajian baris \textit{argument assumption}-nya dibuat berstruktur dan valid runut tanpa menumpuk celah deduksi \textit{jump logic/fallacy} iteratif kalkulatifnya.
\end{asesmen}

\begin{checklist}
  \item Tak bingung menetapkan formula parameter manipulasi genap komputasional $2k$ vs substitusi manipulasi algoritmis logik var variabel ganjil bernilai fungsional rasional $2k+1$.
  \item Secara mahir dan fasih membedah argumen dengan pemisalan \textit{Contrapositive} bilamana *Direct Proof*-nya menyulitkan dan \textit{Stuck} faktorisasi persamaannya.
\end{checklist}

\end{document}
